% Generated draft: requires mathematical and source review.
\documentclass[11pt]{article}
\usepackage[margin=25mm]{geometry}
\usepackage{fontspec}
\setmainfont{Latin Modern Roman}
\usepackage{amsmath,amssymb,mathtools,unicode-math}
\setmathfont{Latin Modern Math}
\usepackage{enumitem,needspace,xurl,hyperref}
\setlength{\parindent}{0pt}
\setlength{\parskip}{5pt}
\setlength{\emergencystretch}{3em}
\newcommand{\N}{\mathbb{N}}
\newcommand{\Z}{\mathbb{Z}}
\newcommand{\Q}{\mathbb{Q}}
\newcommand{\R}{\mathbb{R}}
\newcommand{\C}{\mathbb{C}}
\newcommand{\F}{\mathbb{F}}
\newcommand{\A}{\mathbb{A}}
\newcommand{\PP}{\mathbb{P}}
\newcommand{\E}{\mathbb{E}}
\newcommand{\eps}{\varepsilon}
\newcommand{\dd}{\,\mathrm{d}}
\begin{document}
\section*{UnsolvedMath ID: 700718}
\textbf{Source problem code:} catalogue entry
\subsection{Definitions and mathematical statement}
In the standard finite-group convention, a \emph{Fitting class} is a nonempty class of finite groups, invariant under isomorphism, closed under taking normal subgroups, and closed under normal products: if \(N_1,N_2\trianglelefteq G\) belong to the class, then \(N_1N_2\) belongs to the class. A \emph{formation} is a nonempty class of finite groups closed under homomorphic images and finite subdirect products. Write \(\mathcal G_{\mathrm{fin}}\) and \(\mathcal S_{\mathrm{fin}}\) for the classes of all finite groups and all finite soluble groups, respectively.

The record does not supply the definition of a \emph{local} Fitting class, the cited \emph{product} of Fitting classes, or the meaning of a \emph{non-trivial} product. These notions must be recovered from the cited source before this can become a self-contained statement. In particular, no identification of that product with an ordinary extension product is asserted. The intended ambient universe for part (b), including whether it is restricted to finite soluble groups, also requires confirmation.

The following notation records the quantifiers without resolving those omissions. Let \(\mathfrak F\) denote the collection of Fitting classes in the intended ambient universe, and let \(\operatorname{Loc}(\mathcal C)\) mean that \(\mathcal C\) is local in the source's sense. For \(\mathcal C\in\mathfrak F\), let \(\mathscr D(\mathcal C)\) denote its admissible product decompositions under the cited product, and let \(\mathscr D_{\mathrm{nt}}(\mathcal C)\) denote those considered non-trivial by the source. A decomposition \(d\) has factors \((\mathcal C_i(d))_{i\in I(d)}\). The admissible index sets, product convention, and non-triviality condition are unresolved source parameters.

\textbf{Part (b), the remaining question in the supplied record.} Does there exist \(\mathcal C\in\mathfrak F\) such that
\[
\operatorname{Loc}(\mathcal C),\qquad
\mathscr D_{\mathrm{nt}}(\mathcal C)\ne\varnothing,
\qquad
\forall d\in\mathscr D_{\mathrm{nt}}(\mathcal C)\;
\forall i\in I(d),\quad
\neg\operatorname{Loc}(\mathcal C_i(d))?
\]
Thus at least one non-trivial product decomposition must exist, and every factor in every such decomposition must be non-local. Part (b) does not explicitly repeat the exclusions or the non-formation condition appearing in part (a); they have not been added here.

\textbf{Part (a), preserved as archived context.} Does there exist an admissible product of Fitting classes whose resulting class \(\mathcal C\) satisfies
\[
\operatorname{Loc}(\mathcal C),\qquad
\mathcal C\ne\mathcal G_{\mathrm{fin}},\qquad
\mathcal C\ne\mathcal S_{\mathrm{fin}},
\]
and whose every factor is both non-local and not a formation? In the notation above, this asks for \(d\in\mathscr D(\mathcal C)\) such that
\[
\forall i\in I(d),\qquad
\neg\operatorname{Loc}(\mathcal C_i(d))
\quad\text{and}\quad
\mathcal C_i(d)\text{ is not a formation}.
\]
The supplied archive reports an affirmative answer to part (a), attributed to N. T. Vorob'ev and A. N. Skiba (1995). This is a supplied historical report, not an independently verified result in this draft.
\subsection{Short English statement}
Part (b) asks whether a local Fitting class can have a non-trivial product decomposition while every factor in every such decomposition is non-local. Archived part (a) asks for a local product, different from both the class of all finite groups and the class of all finite soluble groups, whose factors are non-local and are not formations; the supplied archive reports an affirmative answer. The meanings of locality, product, non-triviality, and the ambient universe require source confirmation.
\subsection{Sources}
\begin{itemize}
\item Local registry record: permanent UnsolvedMath ID \(700718\), exact source problem code \texttt{catalogue entry}, title \emph{Kourovka Notebook Problem 11.25}; supplied upstream revision \texttt{d371ee917ac9647fbb06cd6fc529deec6f3dda82}.
\item \emph{The Kourovka Notebook (Unsolved Problems in Group Theory)}, 11th issue (1990), Problem 11.25, proposed by N. T. Vorob'\"ev. The supplied record describes an update in the 21st issue (2026), \texttt{arXiv:1401.0300v48}. Its collection description names E. I. Khukhro and V. D. Mazurov as editors.
\item Product-definition reference supplied with part (b): N. T. Vorob'\"ev, \emph{Math. Notes}, volume \textbf{43}, no. 2 (1988), pp. 91--94. The archived part (a) instead gives N. T. Vorob'ev, \emph{Math. Notes}, volume \textbf{43}, nos. 1--2 (1988), pp. 91--94. No article title or link is supplied; the issue discrepancy remains unresolved.
\item Supplied Archive of Solved Problems attribution for the affirmative answer associated with part (a): N. T. Vorob'ev and A. N. Skiba, \emph{Problems in Algebra}, volume \textbf{8}, Gomel', 1995, pp. 55--58 (Russian). No article title or link is supplied.
\item Supplied status curation checked on \texttt{2026-10-16}: canonical status \texttt{partially\_solved}, evidence label \texttt{published}; the accompanying note places solved parts in the archive and remaining parts in the main section. This date and assessment are supplied dated curation, not independently verified source or literature checking.
\end{itemize}
\end{document}
